\section{Proofs of the state and geometry results}\label{app:r24proofs}
\subsection{Proof of Proposition~\ref{prop:r24realization}}
The conditional-price identity gives $X_i(t)=F_i(t,Y_t)$ and
\[
 dX_i=\{(.02+.06p_i)X_i-c_i\}dt+.2p_iX_i\,dW^x.
\]
Multiply by constant $\lambda_i$ and sum. On \eqref{eq:r24manifold}, this is the original wealth equation with \eqref{eq:r24policy}. All virtual preference processes share the same shock and $\sum_i\lambda_i=1$, so their weighted sum solves $du=\theta^\lambda dt+.05dW^u$. Pathwise uniqueness on each slab and continuity of prices at slab dates preserve both equalities until aggregate stopping. Consumption and drift are convex combinations within their admissible intervals. Portfolio exposure is a wealth-weighted convex combination because $F_i>0$ and $x=\sum_i\lambda_iF_i$. The global bounds of Theorem~\ref{thm:r21finance} give admissibility. Frozen coefficients, constant weights, and current $(t,Y_t)$ determine prices and actions. Preference differences are initial differences plus integrated deterministic drift differences, requiring no additional stochastic coordinate.\qed

\subsection{Proof of Theorem~\ref{thm:r24value}}
Fix deterministic initial memory. Every augmented progressively measurable control is an original progressively measurable control: both use the same filtration, and memory is constructed causally from it. The economic coordinates have identical dynamics and payoff under that control. Hence the augmented supremum is at most $V_k(t,u,x)$. Conversely, every original admissible control is available by ignoring memory, giving the reverse inequality. This compares admissible classes and does not require that a neural family attain the supremum.

Covariances follow from diffusion coefficients and correlations. In particular,
\[
 d\langle u,Y\rangle_t=.05(-.3Y_t)(-.25)dt=.00375Y_tdt,\qquad
 d\langle x,Y\rangle_t=-.06p_tx_tY_tdt.
\]
Also $d\langle Y\rangle_t=.09Y_t^2dt$. The factor $1/2$ applies to diagonal Hessian terms, not mixed terms written once. There are no derivatives in constant weights. It\^o's formula gives \eqref{eq:r24generator}.

Conditioning at an intermediate time settles paths already stopped and uses admissible concatenation on surviving paths, giving \eqref{eq:r24semigroup}. If $w\leq\widetilde w$, each admissible control has weakly larger continuation payoff under $\widetilde w$; taking suprema preserves order. Conditioning and concatenating twice gives the semigroup identity. A classical first-order time expansion, when justified, gives \eqref{eq:r24HJB}. Memory derivatives of the unrestricted value vanish, reducing it to the original Bellman equation. The value comparison itself needs no differentiability.\qed

\subsection{Proof of Theorem~\ref{thm:r24transport}}
The chain rule gives $\nabla_\vartheta L=J_n'g_n$. Iterate $m_n=\beta_1m_{n-1}+(1-\beta_1)J_n'g_n$ from $m_0=0$ and divide by $1-\beta_1^n$ to obtain the historical expansion. Adam uses $\delta\vartheta_n=-\eta_nD_n^{-1}\widehat m_n$. The integral remainder is
\[
 f(\vartheta_n+\delta\vartheta_n)-f(\vartheta_n)-J_n\delta\vartheta_n
 =\int_0^1\{Df(\vartheta_n+s\delta\vartheta_n)-J_n\}\delta\vartheta_n ds.
\]
Its norm is at most $H\|\delta\vartheta_n\|^2/2$. Substitute the update to obtain the two displayed conclusions. For any $v$, $v'JD_n^{-1}J'v=\|D_n^{-1/2}J'v\|^2\geq0$; rank cannot exceed the output dimension. On an output-gradient subspace with smallest singular value $s_*>0$ for $J'$, $\|J'g\|\geq s_*\|g\|$. Without that lower bound, small parameter gradient can coexist with large output gradient. Historical $J_r$ cannot generally be replaced by $J_n$; even equality of Jacobians leaves an average of historical gradients, not necessarily the current gradient.\qed

\subsection{Proof of Proposition~\ref{prop:r24charts}}
The chain rule yields $D\widetilde f=JS$. A gradient step is $\delta\varphi=-\eta S'J'g$, so $J\delta\vartheta=JS\delta\varphi=-\eta JSS'J'g$. Local Lipschitz continuity of the Jacobian bounds the nonlinear remainder by a constant times the squared parameter step. If $J$ has full row rank, $JJ'$ is invertible and the stated minimum-norm right inverse satisfies $J\delta\vartheta=-\eta d$ exactly in the linearization. Other full-rank parameter charts recover that prescribed first-order output direction, although step norm and nonlinear remainder can differ.\qed

\subsection{Proof of Theorem~\ref{thm:r24gradient}}
At the approximate root, $|P(y,a_N)-B_0|=|P(y,a_N)-P_N(y,a_N)|\leq e_P$. Monotonicity and the mean-value theorem imply $|a_N-a|\leq e_P/d_*$. A fixed-offset derivative error $e_X$ and offset-Lipschitz constant $H_X$ imply a root-evaluated error $E_X=e_X+H_X\delta$ by the triangle inequality. Implicit differentiation gives
\[
 \nabla_y F(y,a(y))=F_y-F_aP_y/P_a,
\]
and similarly for the numerical functions. Abbreviate derivatives by $A,B,C,D$ and numerical versions by $A_N,B_N,C_N,D_N$. Add and subtract $BC_N/D_N$ and $BC/D_N$ to obtain
\[
 \|B_NC_N/D_N-BC/D\|
 \leq {|B_N-B|\|C_N\|\over d_*}
 +{|B|\|C_N-C\|\over d_*}
 +{|B|\|C\||D_N-D|\over d_*^2}.
\]
Adding the $A$ error proves \eqref{eq:r24gradienterror}. Every derivative bound is a hypothesis, not an inference from an ordinary two-rule difference or a payoff-value error bound.\qed

\subsection{Proof of Proposition~\ref{prop:r24cover}}
The retained common-shock mixture argument gives $J(0,z;\pi^\lambda)\geq\sum_i\lambda_iL_i-\zeta$ under its concavity, integrability and stopping hypotheses. Convexity of the explicit upper gives $V(0,z)\leq\overline V(z)\leq\sum_i\lambda_i\overline V(z_i)$. Subtract and bound the convex combination of vertex gaps by its largest term. Maximizing over a complete finite cover proves the union bound. At a shared face a deterministic initialization rule chooses a cell; either neighboring certified construction supplies its corresponding bound. This is not repeated switching of controllers during a path. The result is a regret bound. Improved vertex payoffs alone do not prove uniform mixture payoff improvement, because the old and new mixtures can have different Jensen gains; the retained separate mixture-gain theorem addresses that stronger comparison.\qed

\subsection{Proof of Proposition~\ref{prop:r24reference}}
The function $c\mapsto-1/c$ is strictly concave for $c>0$. At a fixed average $\bar c$, Jensen's inequality gives $\int_0^1-1/c_tdt\leq-1/\bar c$, with equality only for constant consumption almost everywhere. Terminal wealth is $1.25-\bar c$. Thus maximize $h(c)=-1/c+\log(1.25-c)$, whose second derivative is $-2/c^3-1/(1.25-c)^2<0$. Its first-order condition gives $c^2+c=1.25$ and positive root $(\sqrt6-1)/2$. Direct substitution verifies interior control bounds and reserve above $.5$. Exact integration on each slab gives the stated finite-slab payoff. Constant output bias represents the optimum in both continuous parameter families, so there is no policy-class loss. The optimum and candidate expressions require rational operations, a square root and a logarithm, enclosed by the directed backend. Their interval difference encloses the optimization gap rather than estimating it statistically.\qed
